% =========================================================
% Data Description — UCI Online Retail Dataset
% For inclusion in the revised manuscript (Section 2 or 6)
% =========================================================

\subsection{UCI Online Retail Dataset}\label{sec:uci_data}

To complement the ASOS dataset analyzed in the original submission,
we consider the UCI Online Retail dataset \citep{chen2012data},
a publicly available transactional record from a UK-based non-store
online retailer covering December~2010 through December~2011.
The raw data comprises 541{,}909 transaction line items across
4{,}339 unique customers over 374~calendar days.

\paragraph{Preprocessing.}
We retain only transactions with a valid customer identifier and
positive quantity, yielding 397{,}924 rows.
For each customer~$n$ and calendar day~$d$, we define the
\emph{trigger count} $A_{d,n}$ as the number of distinct transaction
line items recorded for that customer on that day.
This produces 16{,}766 customer--day pairs.

\paragraph{User engagement characteristics.}
The dataset exhibits the heavy-tailed engagement pattern
that motivates our methodology:
35.7\% of customers are active on exactly one day,
67.4\% on three or fewer days,
while 0.6\% are active on more than 30~days.
The median total trigger count per customer is~41,
but the distribution has a long right tail
(maximum 7{,}847 triggers; standard deviation 228.8).
These features are consistent with the power-law behavior
discussed in Section~\ref{sec:setting}
and documented broadly in online user activity data
\citep{clauset2009power}.

\paragraph{Experiment construction.}
We partition the 374-day observation period into
13~non-overlapping windows of $D = 28$~days each.
Each window is treated as an independent ``experiment.''
Within each window, the first $D_0 = 7$~days serve as the
pilot period and the remaining $D_1 = 21$~days as the follow-up.
Table~\ref{tab:uci_experiments} summarizes the resulting experiments.
The windows span a range of traffic levels:
from a low-activity period in early January
(Experiment~1, $N_{D_0} = 34$)
to the holiday season
(Experiment~12, $N_{D_0} = 474$),
providing a natural test of model robustness across
varying data sparsity.

\begin{table}[t]
\centering
\caption{Summary of the 13 experiment windows constructed from the
UCI Online Retail dataset.
$N_{D_0}$: users active in the 7-day pilot;
$N_{\mathrm{total}}$: users active in the full 28-day window;
$U_{D_0}^{(D_1)}$: new users in the 21-day follow-up;
$T_{D_0}^{(D_1)}$: total triggers in the follow-up.}
\label{tab:uci_experiments}
\small
\begin{tabular}{rrrrrr}
\toprule
Exp & Days & $N_{D_0}$ & $N_{\mathrm{total}}$ & $U_{D_0}^{(D_1)}$ & $T_{D_0}^{(D_1)}$ \\
\midrule
0  & 0--28   & 423 & 885   & 462   & 15{,}352 \\
1  & 28--56  & 34  & 606   & 572   & 15{,}827 \\
2  & 56--84  & 249 & 751   & 502   & 14{,}027 \\
3  & 84--112 & 264 & 846   & 582   & 17{,}153 \\
4  & 112--140 & 278 & 960  & 682   & 19{,}381 \\
5  & 140--168 & 170 & 854  & 684   & 18{,}552 \\
6  & 168--196 & 348 & 940  & 592   & 17{,}468 \\
7  & 196--224 & 285 & 861  & 576   & 16{,}333 \\
8  & 224--252 & 277 & 909  & 632   & 19{,}043 \\
9  & 252--280 & 242 & 878  & 636   & 19{,}545 \\
10 & 280--308 & 331 & 1{,}187 & 856 & 30{,}014 \\
11 & 308--336 & 418 & 1{,}307 & 889 & 33{,}518 \\
12 & 336--364 & 474 & 1{,}589 & 1{,}115 & 47{,}036 \\
\bottomrule
\end{tabular}
\end{table}

\paragraph{Comparison with the ASOS dataset.}
Unlike the ASOS data, which provides only the number of distinct
users first triggering on each day, the UCI dataset provides
per-customer daily trigger counts.
This richer structure allows us to fit and evaluate the
\nbmodel model in addition to the \bemodel and \tgmodel models,
and to assess predictions of total future triggers
$T_{D_0}^{(D_1)}$ on public data for the first time.
